% !TeX root = ../report.tex
\documentclass[../report.tex]{subfiles}
\begin{document}

\section{原系统根轨迹}
对于原系统的开环传递函数：
\begin{equation}
    G(s)H(s) = \frac{K(s+20)}{s(s^2 + 24s + 104)}
\end{equation}

本节第一部分绘制根轨迹图并计算相关关键点（渐近线，分离/汇合点，出射/入射角，和虚轴交点），第二部分分析系统稳定性，
第三部分分析阻尼比 $\xi$ 和自然频率 $\omega_n$ 在根轨迹上的几何意义，第四部分根据超调量选择合适的增益 $K$ 。

%%%%%%%%%%%%%%%%%%%%%%%%%%%%%%%%%%%%%%%%%%%%%%%%%%%%%%%%%%%%%%%%%%%%%%%%%%%%%%%%%%%%%%%%%%%%%%%%%%%%%%%%%%%%
\subsection{根轨迹的绘制}

\subsubsection{根轨迹的绘制：}
利用 MATLAB 可以较为方便地绘制根轨迹图，代码如下：

\begin{lstlisting}[language=Matlab]
    s= tf('s');
    gs = (s+20)/(s*(s*s + 24*s +104));
    r = rlocus(gs);
\end{lstlisting}

得到的根轨迹图像如图\ref{fig:ques1_rlocus}所示。
\begin{figure}[H]
    \centering
    \includegraphics[width=0.8\textwidth]{./ques1/rlocus.png}
    \caption{根轨迹图}
    \label{fig:ques1_rlocus}
\end{figure}
可以看到，本系统有3个开环极点，1个开环零点。根轨迹从3个开环极点出发，最终汇聚到1个开环零点和2个无穷远处的渐近线上。

\subsubsection{渐近线的绘制：}
根轨迹的渐近线倾斜角公式为：
\begin{equation}
    \theta_a = \frac{(2k+1)\pi}{n-m}, k=0,1,2,\cdots
\end{equation}
其中，$n$ 为系统开环传递函数的极点个数，$m$ 为系统开环传递函数的零点个数。对于本系统，$n=3, m=1$，因此渐近线的角度为：
\begin{equation}
    \theta_a = \frac{(2k+1)\pi}{3-1} = \frac{(2k+1)\pi}{2}, k=0,1,2,\cdots
\end{equation}
根轨迹的渐近线与实轴交点公式为：
\begin{equation}
    \sigma_a = \frac{\sum_{i=1}^{n}p_i - \sum_{j=1}^{m}z_j}{n-m}
\end{equation}
经过计算可以得到，$\sigma_a = -2$，因此渐近线与实轴交点为 $-2$。

\subsubsection{分离/汇合点的绘制：}
根轨迹分离/汇合点公式为：
\begin{equation}
    \sum_{i=1}^{n}\frac{1}{s - p_i} - \sum_{j=1}^{m}\frac{1}{s - z_j} = 0
\end{equation}
经过计算可以得到，分离/汇合点为 $s = -2.81, K = 7,27$。

\subsubsection{出射/入射角的求解：}
出射/入射角可以使用以下公式计算：
\begin{align}
    \theta &= 180^\circ + \angle(s - z_1) + \angle(s - z_2) + \cdots + \angle(s - z_m) - \angle(s - p_1) - \angle(s - p_2) - \cdots - \angle(s - p_n)
\\
    \varphi &= 180^\circ + \angle(s - p_1) + \angle(s - p_2) + \cdots + \angle(s - p_n) - \angle(s - z_1) - \angle(s - z_2) - \cdots - \angle(s - z_m)
\end{align}
然而，由于本系统的开环零极点均分布在实轴上，因此出射/入射角均为 $0^\circ$或$180^\circ$， 可以直接由图得到，并不需要使用此公式。

\subsubsection{和虚轴交点的计算：}
和虚轴交点可以利用特征方程 $D(s) = 0$ 来求解。以本系统为例，对$1+G(s)H(s)$进行化简，得到特征方程为：
\begin{equation}
    D(s) = s^3 + 24s^2 + (104+K)s + 20K = 0
\end{equation}
假定有根分布在虚轴上，也就是 $s = j\omega$，代入特征方程得到：
\begin{equation}
    (j\omega)^3 + 24(j\omega)^2 + (104+K)(j\omega) + 20K = 0
\end{equation}
将其分为实部和虚部，得到：
\begin{equation}
    \begin{cases}
        20K - 24\omega^2 = 0 \\
        -\omega^3 + (104+K)\omega = 0
    \end{cases}
\end{equation}
求解得到：$\omega = 0, K = 0$ 或 $\omega = j23.37, K = -624.5$（$K<0, \omega \notin \mathbb{R} $， 舍去），
因此本系统除原点外没有和虚轴交点。

综上，本部分得到结论：
\begin{itemize}
    \item 渐近线与实轴交点为 $-2$，渐近线角度为 $90^\circ, 270^\circ$；
    \item 分离/汇合点为 $s = -2.81, K = 7.27$，（见图\ref{fig:ques1_rlocus_ori}）；
    \item 出射/入射角均为 $0^\circ$ 或 $180^\circ$；
    \item 除原点外没有和虚轴交点。
\end{itemize}

\begin{figure}[H]
    \centering
    \includegraphics[width=0.8\textwidth]{./ques1/rlocus_ori.png}
    \caption{根轨迹图}
    \label{fig:ques1_rlocus_ori}
\end{figure}

%%%%%%%%%%%%%%%%%%%%%%%%%%%%%%%%%%%%%%%%%%%%%%%%%%%%%%%%%%%%%%%%%%%%%%%%%%%%%%%%%%%%%%%%%%%%%%%%%%%%%%%%%%%%%
\subsection{系统稳定性分析}
根轨迹的稳定性分析可以通过观察根轨迹图来进行。稳定性判据为：当系统的闭环极点全部位于左半平面时，系统是稳定的；
当系统的闭环极点有一个或多个位于右半平面时，系统是不稳定的；当系统的闭环极点有一个或多个位于虚轴上时，系统是临界稳定的。

对于本系统，根轨迹从三个开环极点出发，其中两个极点位于左半平面，一个极点位于原点。随着增益 $K$ 
的增加，根轨迹会向右移动，并最终汇聚到开环零点和无穷远处的渐近线上。

因而可以得到结论：
\begin{itemize}
    \item 当 $K=0$ 时，系统的闭环极点为 $s = 0, -12 \pm j8$，其中一个极点位于原点，系统是临界稳定的；
    \item 当 $0 < K \leq  7.27$ 时，系统的闭环极点全部位于左半实轴上，系统是稳定的，并且单调收敛；
    \item 当 $K > 7.27$ 时，系统闭环极点全部位于左半平面上，系统是稳定的，并且由一对共轭极点，系统震荡收敛。
\end{itemize}

%%%%%%%%%%%%%%%%%%%%%%%%%%%%%%%%%%%%%%%%%%%%%%%%%%%%%%%%%%%%%%%%%%%%%%%%%%%%%%%%%%%%%%%%%%%%%%%%%%%%%%%%%%%%%%
\subsection{阻尼比 $\xi$ 和自然频率 $\omega_n$ 的几何意义}\label{sec:wxi}

对于一对共轭极点 $s = \sigma \pm j\omega_d$，其阻尼比 $\xi$ 和自然频率 $\omega_n$ 的几何意义如下：
\begin{itemize}
    \item 阻尼比 $\xi$：表示系统的阻尼程度，定义为 $\xi = -\frac{\sigma}{\sqrt{\sigma^2 + \omega_d^2}}$。
    在根轨迹图中，阻尼比 $\xi$ 可以通过极点到虚轴的距离与极点到原点的距离之比来表示。具体来说，阻尼比 $\xi$ 越大，
    极点越靠近实轴，系统的响应越快且振荡越小；反之，阻尼比 $\xi$ 越小，极点越靠近虚轴，系统的响应越慢且振荡越大。
    \item 自然频率 $\omega_n$：表示系统的固有振荡频率，定义为 $\omega_n = \sqrt{\sigma^2 + \omega_d^2}$。
    在根轨迹图中，自然频率 $\omega_n$ 可以通过极点到原点的距离来表示。具体来说，自然频率 $\omega_n$ 越大，
    系统的响应速度越快；反之，自然频率 $\omega_n$ 越小，系统的响应速度越慢。
\end{itemize}

下面不妨用二阶系统来举例说明，假设系统的闭环传递函数为：
\begin{equation}
    T(s) = \frac{\omega_n^2}{s^2 + 2\xi\omega_n s + \omega_n^2}
\end{equation}

首先固定 $\omega_n$，改变 $\xi$ 的值，观察系统的闭环极点在根轨迹图上的位置变化。可以发现，当 $\xi$ 增大时，闭环极点会向实轴靠近，系统的响应速度加快，振荡减小；当 $\xi$ 减小时，闭环极点会向虚轴靠近，系统的响应速度减慢，振荡增大。
然后固定 $\xi$，改变 $\omega_n$ 的值，观察系统的闭环极点在根轨迹图上的位置变化。可以发现，当 $\omega_n$ 增大时，闭环极点远离原点，系统的响应速度加快；当 $\omega_n$ 减小时，闭环极点会靠近原点，系统的响应速度减慢。
具体结果如图\ref{fig:ques1_xi_wn}所示。
\begin{figure}
    \centering
    \includegraphics[width=0.8\textwidth]{./ques1/exp.png}
    \caption{二阶系统中的阻尼比 $\xi$ 和自然频率 $\omega_n$ }
    \label{fig:ques1_xi_wn}
\end{figure}

高阶系统的闭环极点可能有多个共轭极点对，可能需要考虑每一对共轭极点的阻尼比和自然频率。当
系统多对共轭极点距离虚轴距离差别较大时，可以只考虑系统主极点，从而将高阶系统近似为二阶系统。故而，
阻尼比 $\xi$ 和自然频率 $\omega_n$ 的几何意义在高阶系统中仍然适用。

%%%%%%%%%%%%%%%%%%%%%%%%%%%%%%%%%%%%%%%%%%%%%%%%%%%%%%%%%%%%%%%%%%%%%%%%%%%%%%%%%%%%%%%%%%%%%%%%%%%%%%%%%%%%%%
\subsection{根据超调量选择合适的增益 $K$}
由图\ref{fig:ques1_rlocus_ori}可以看到，系统由三个极点，但是最左侧极点距离虚轴较远，因此可以仅考虑右侧两极点，
将系统近似为二阶系统。当$K \leq 7.27$时，系统保持稳定而且单调收敛，故而系统零超调。当$K > 7.27$时，系统近似
为欠阻尼二阶系统如下：
\begin{equation}
    T(s) = \frac{\omega_n^2}{s^2 + 2\xi\omega_n s + \omega_n^2}
\end{equation}
其闭环传递函数是：
\begin{equation}
    T(s) = \frac{\omega_n^2}{s^2 + 2\xi\omega_n s + \omega_n^2}
\end{equation}
输入阶跃响应后，系统的响应是：
\begin{equation}
    c(t) = 1 - \frac{1}{\sqrt{1-\xi^2}}e^{-\xi\omega_n t}\sin(\omega_d t + \phi)
\end{equation}
其中，$\omega_d = \omega_n\sqrt{1-\xi^2}$， $\phi = \arccos(\xi)$。

对$c(t)$求导，得到：
\begin{equation}
    \frac{dc(t)}{dt} = \frac{1}{\sqrt{1-\xi^2}}e^{-\xi\omega_n t}(\xi\omega_n\sin(\omega_d t + \phi) - \omega_d\cos(\omega_d t + \phi)) 
\end{equation}
令 $\frac{dc(t)}{dt} = 0$，得到：
\begin{equation}
    \xi\omega_n\sin(\omega_d t + \phi) - \omega_d\cos(\omega_d t + \phi) = 0
\end{equation}
从而有$c(t)$的极值点为：
\begin{equation}
    t_p = \frac{\pi}{\omega_d} = \frac{\pi}{\omega_n\sqrt{1-\xi^2}}
\end{equation}
将$t_p$代入$c(t)$，得到系统的最大超调量为：
\begin{align}
    M_p = c(t_p) - 1 = e^{-\frac{\pi\xi}{\sqrt{1-\xi^2}}}
\end{align}
令 $M_p = 0.1$，得到 $\xi = 0.5912$。
结合图像\ref{fig:ques1_rlocus_K}，可以得到 $\xi = 0.5912$ 对应的增益 $K = 20.3$，在时域上验证，符合预期（见图\ref{fig:ques1_time_K}），超调量为 $\sigma \% = 9.95\%$。
\begin{figure}[H]
    \centering
    \begin{subfigure}[b]{0.4\textwidth}
        \centering
        \includegraphics[width=\linewidth]{./ques1/rlocus10.png}
        \caption{根轨迹图}
        \label{fig:ques1_rlocus_K}
    \end{subfigure}
    \begin{subfigure}[b]{0.4\textwidth}
        \centering
        \includegraphics[width=\linewidth]{./ques1/time10.png}
        \caption{时域响应图}
        \label{fig:ques1_time_K}
    \end{subfigure}
    \caption{增益 $K$ 对系统超调量的影响}
\end{figure}

故而，当系统的超调量要求为 $10\%$ 时，选择增益 $0< K < 20.3$ 可以满足要求。


\end{document}